\chapter{Morphosyntax: Structural Animacy Systems}

\section{The Active-Stative Alignment Framework}
The morphological architecture of the PED reconstruction is defined by an Active-Stative alignment. In this framework, nouns are categorized according to their inherent capacity for agency and volition. This system is hypothesized to be the ancestral state from which the Nominative-Accusative and Ergative-Absolutive systems of Eurasia later emerged.

\subsection{Agential Categorization (*-S)}
Beings perceived as capable of independent motion and action—including humans, animals, and natural forces—were marked with the \textit{*-S} suffix. 
\begin{itemize}
    \item \textbf{Vedic Reflex:} Corresponds to the Masculine and Feminine nominative singular \textit{*-s}.
    \item \textbf{Old Tamil Reflex:} Preserved in the 'Rational' (uyartiṇai) gender classifications.
\end{itemize}

\subsection{Resultative Categorization (*-N)}
Passive objects, experiences, and resultative nouns were marked with the \textit{*-N} suffix.
\begin{itemize}
    \item \textbf{Vedic Reflex:} Corresponds to the Neuter nominative/accusative \textit{*-m}.
    \item \textbf{Old Tamil Reflex:} Preserved in the 'Irrational' (akṟiṇai) neuter singular markers.
\end{itemize}

\section{Noun Declension Paradigm}
The PED noun was a modular unit carrying functional suffixes. The following table maps the primary ancestral markers to their attested reflexes in the earliest daughter strata.

\begin{longtable}{@{}llll@{}}
\toprule
\textbf{Functional Role} & \textbf{PED Suffix} & \textbf{Vedic Reflex} & \textbf{Tamil Reflex} \\ \midrule
\endfirsthead
\toprule
\textbf{Functional Role} & \textbf{PED Suffix} & \textbf{Vedic Reflex} & \textbf{Tamil Reflex} \\ \midrule
\endhead
\bottomrule
\endfoot
\bottomrule
\lastfoot
Agential & *-S & Nom. *-s & Rational Marker \\
Resultative & *-N & Neut. *-m & Accusative *-n \\
Directional & *-K & (and) *-kʷe & Dative *-ku \\
Locative & *-R & Loc. *-er & Locative *-il \\
Instrumental & *-P & Plural *-bhi & Increment *-p- \\
Reflexive & *-v- & (Middle Voice) & koḷ- (Take) \\
\end{longtable}

\section{Deictic Agreement and Proximity}
Agreement markers reflected the spatial relationship between the speaker and the referent, utilizing a tripartite deictic system (\textit{i/u/a}). We identify the interaction of agential markers with these deictic nodes.

\begin{table}[h]
\centering
\begin{tabular}{@{}llll@{}}
\toprule
\textbf{Agential Status} & \textbf{Proximate (*i)} & \textbf{Medial (*u)} & \textbf{Distal (*a)} \\ \midrule
Active Self (1st) & *eK & (n/a) & (n/a) \\
Active Other (2nd) & (n/a) & *tV & (n/a) \\
Stative Other (3rd) & *i-N & *u-N & *a-N \\
Active Other (3rd) & *i-S & *u-S & *a-S \\ \bottomrule
\end{tabular}
\caption{The Hyper-Detailed Deictic Agreement Matrix.}
\end{table}

\section{Structural Fossils: The *K-N* Osteological Link}
A primary evidence for the PED substrate is the existence of 'broken' paradigms that defy coincidental invention. We identify the \textit{*K-N} cluster governing the concept of 'joint', 'angle', or 'bend'.
\begin{center}
\textit{*K-N} $\rightarrow$ Vedic \textit{akṣi} (joint/eye), Old Tamil \textit{kaṇ} (joint/eye), Old Tibetan \textit{m-kun} (knee).
\end{center}
The arbitrary mapping of 'joint' to 'eye' (the articulating point of the face) across all three families provides a diagnostic structural invariant.

\section{Transitive Consonant Hardening}
PED utilized a phonological mechanism to indicate the transition from a stative to a causative/transitive state. This involved the hardening or gemination of the root-initial consonant.
\begin{itemize}
    \item \textbf{Southern Node (Dravidian):} Preserved as the gemination rule (e.g., \textit{*āṭu} "to dance" vs. \textit{*āṭṭu} "to cause to dance").
    \item \textbf{Western Node (Indo-European):} Preserved as the \textit{s-mobile} prefix, where the addition of agential friction modified the root onset (e.g., \textit{*(s)teg-}).
\end{itemize}
This 'Causal Engine' provides the developmental link between the primitive PED root and the complex verbal systems of the daughter languages.
